
% ============================================================
% UNIDAD 1. NÚMEROS COMPLEJOS
% PARTE I: FUNDAMENTOS ALGEBRAICOS
% Versión ampliada, formal y con ejemplos para guía docente
% ============================================================

\section{Números complejos}

\subsection{Motivación: ampliación de los sistemas numéricos}

Los sistemas numéricos que se utilizan en matemáticas pueden entenderse como
extensiones sucesivas construidas para resolver ecuaciones u operaciones que
no son posibles dentro de un sistema previo.

Por ejemplo, la ecuación
\[
x+3=1
\]
no posee solución en $\mathbb{N}$, pero sí en $\mathbb{Z}$. De manera análoga,
\[
2x=1
\]
no posee solución en $\mathbb{Z}$, pero sí en $\mathbb{Q}$. Finalmente,
\[
x^2=2
\]
no posee solución en $\mathbb{Q}$, pero sí en $\mathbb{R}$.

Esto conduce a la cadena de inclusiones
\[
\mathbb{N}\subset\mathbb{Z}\subset\mathbb{Q}\subset\mathbb{R}.
\]

Sin embargo, incluso el sistema de los números reales es insuficiente para
resolver todas las ecuaciones polinómicas.

Consideremos
\[
x^2+1=0.
\]
La ecuación equivale a
\[
x^2=-1.
\]
Pero para todo $x\in\mathbb{R}$,
\[
x^2\geq 0.
\]
Por tanto,
\[
x^2=-1
\]
no posee solución en $\mathbb{R}$.

La construcción de los números complejos responde precisamente a la necesidad
de extender $\mathbb{R}$ de manera que esta ecuación, y más generalmente todo
polinomio no constante, pueda estudiarse dentro de un sistema algebraicamente
más completo.

\textbf{Observación docente.}
Es importante enfatizar que una extensión numérica no elimina el sistema
anterior: lo contiene como subconjunto. El objetivo es ampliar el universo de
objetos sin perder las reglas algebraicas ya conocidas.

\paragraph{Ejemplo 1.}
La ecuación
\[
x+5=2
\]
no tiene solución en $\mathbb{N}$, pero sí en $\mathbb{Z}$:
\[
x=-3.
\]

\paragraph{Ejemplo 2.}
La ecuación
\[
3x=1
\]
no tiene solución en $\mathbb{Z}$, pero sí en $\mathbb{Q}$:
\[
x=\frac13.
\]

\paragraph{Ejemplo 3.}
La ecuación
\[
x^2=3
\]
no tiene solución racional, pero sí real:
\[
x=\pm\sqrt3.
\]

\paragraph{Ejemplo 4.}
La ecuación
\[
x^2+4=0
\]
equivale a
\[
x^2=-4.
\]
No posee solución real, pero en el sistema complejo tendrá soluciones
\[
x=\pm 2i.
\]

\paragraph{Ejemplo 5.}
La ecuación
\[
x^2-2x+5=0
\]
tiene discriminante
\[
\Delta=b^2-4ac=(-2)^2-4(1)(5)=4-20=-16.
\]
Por tanto, no tiene raíces reales. Más adelante veremos que sus soluciones
complejas son
\[
x=1\pm 2i.
\]

% ============================================================
\subsection{La unidad imaginaria y sus potencias}
% ============================================================

\textbf{Definición.}
Se introduce un nuevo elemento $i$, denominado \emph{unidad imaginaria}, que
satisface la relación
\[
\boxed{i^2=-1.}
\]

La propiedad anterior es la relación algebraica fundamental de la unidad
imaginaria. A partir de ella se obtienen todas sus potencias.

Tenemos
\[
i^0=1,
\qquad
i^1=i,
\qquad
i^2=-1,
\]
y
\[
i^3=i^2i=-i,
\qquad
i^4=i^2i^2=(-1)(-1)=1.
\]

A partir de $i^4=1$ se sigue que las potencias de $i$ son periódicas de
periodo $4$:
\[
\boxed{i^{n+4}=i^n.}
\]

Más precisamente, si
\[
n=4q+r,
\qquad r\in\{0,1,2,3\},
\]
entonces
\[
i^n=i^r.
\]

\textbf{Proposición.}
Para todo $n\in\mathbb{N}$,
\[
i^n=
\begin{cases}
1, & n\equiv 0 \pmod 4,\\
i, & n\equiv 1 \pmod 4,\\
-1, & n\equiv 2 \pmod 4,\\
-i, & n\equiv 3 \pmod 4.
\end{cases}
\]

\textbf{Demostración.}
Sea $n=4q+r$ con $r\in\{0,1,2,3\}$. Entonces
\[
i^n=i^{4q+r}=(i^4)^qi^r=1^qi^r=i^r.
\]
Basta sustituir los cuatro valores posibles de $r$.
\hfill$\square$

\paragraph{Ejemplo 1.}
\[
i^{18}=i^{4(4)+2}=i^2=-1.
\]

\paragraph{Ejemplo 2.}
\[
i^{35}=i^{4(8)+3}=i^3=-i.
\]

\paragraph{Ejemplo 3.}
\[
i^{100}=i^{4(25)}=1.
\]

\paragraph{Ejemplo 4.}
\[
i^{2026}=i^{4(506)+2}=i^2=-1.
\]

\paragraph{Ejemplo 5.}
Simplifiquemos
\[
i^{17}+i^{28}-i^{39}.
\]
Tenemos
\[
i^{17}=i,\qquad
i^{28}=1,\qquad
i^{39}=i^3=-i.
\]
Por tanto,
\[
i^{17}+i^{28}-i^{39}
=i+1-(-i)
=
\boxed{1+2i}.
\]

% ============================================================
\subsection{Definición de número complejo, parte real e imaginaria}
% ============================================================

\textbf{Definición.}
Un número complejo es una expresión de la forma
\[
\boxed{z=a+bi,}
\]
donde
\[
a,b\in\mathbb{R}
\]
e
\[
i^2=-1.
\]

El conjunto de todos los números complejos se denota por
\[
\boxed{\mathbb{C}}.
\]
Así,
\[
\boxed{
\mathbb{C}
=
\{a+bi:\;a,b\in\mathbb{R}\}.
}
\]

Si
\[
z=a+bi,
\]
se define
\[
\boxed{\operatorname{Re}(z)=a}
\]
y
\[
\boxed{\operatorname{Im}(z)=b.}
\]

Es esencial observar que la parte imaginaria de $z$ es $b$, no $bi$.

La representación
\[
z=a+bi
\]
se denomina \emph{forma binómica}, \emph{forma rectangular} o
\emph{forma cartesiana}.

\textbf{Definición.}
Un número complejo $z=a+bi$ es:
\begin{itemize}
    \item real si $b=0$;
    \item imaginario puro si $a=0$ y $b\neq 0$;
    \item complejo no real si $b\neq 0$.
\end{itemize}

Todo número real puede identificarse con el complejo
\[
a=a+0i.
\]
Por tanto,
\[
\boxed{\mathbb{R}\subset\mathbb{C}.}
\]

\textbf{Proposición de unicidad.}
Si
\[
a+bi=c+di,
\]
con $a,b,c,d\in\mathbb{R}$, entonces
\[
a=c
\qquad\text{y}\qquad
b=d.
\]

Equivalentemente,
\[
\boxed{
a+bi=c+di
\iff
a=c\ \text{y}\ b=d.
}
\]

\paragraph{Ejemplo 1.}
Para
\[
z=4-7i
\]
se tiene
\[
\operatorname{Re}(z)=4,
\qquad
\operatorname{Im}(z)=-7.
\]

\paragraph{Ejemplo 2.}
El número
\[
z=-5
\]
puede escribirse como
\[
z=-5+0i.
\]
Por tanto, es un número real considerado como elemento de $\mathbb{C}$.

\paragraph{Ejemplo 3.}
El número
\[
z=\frac32 i
\]
tiene
\[
\operatorname{Re}(z)=0,
\qquad
\operatorname{Im}(z)=\frac32,
\]
por lo que es imaginario puro.

\paragraph{Ejemplo 4.}
Si
\[
(2x-1)+(y+3)i=7-4i,
\]
la igualdad de números complejos implica
\[
2x-1=7,
\qquad
y+3=-4.
\]
Así,
\[
\boxed{x=4,\qquad y=-7.}
\]

\paragraph{Ejemplo 5.}
Determine $a$ y $b$ si
\[
(a+2)+(3b-1)i=5+8i.
\]
Igualando componentes:
\[
a+2=5,
\qquad
3b-1=8.
\]
Por tanto,
\[
\boxed{a=3,\qquad b=3.}
\]

% ============================================================
\subsection{Construcción formal de $\mathbb{C}$ a partir de $\mathbb{R}^2$}
% ============================================================

La notación $a+bi$ es conveniente, pero los números complejos pueden
construirse rigurosamente sin asumir previamente la existencia de $i$.

Consideremos
\[
\mathbb{R}^2
=
\{(a,b):a,b\in\mathbb{R}\}.
\]

Definimos la suma
\[
(a,b)+(c,d)
=
(a+c,b+d)
\]
y la multiplicación
\[
\boxed{
(a,b)(c,d)
=
(ac-bd,ad+bc).
}
\]

Identificamos cada número real $a$ con el par
\[
(a,0).
\]
De esta manera,
\[
0=(0,0),
\qquad
1=(1,0).
\]

Definimos
\[
i=(0,1).
\]
Entonces
\[
i^2
=
(0,1)(0,1)
=
(-1,0),
\]
y como $(-1,0)$ representa al número real $-1$,
\[
\boxed{i^2=-1.}
\]

Además,
\[
(a,b)
=
(a,0)+(b,0)(0,1),
\]
lo que justifica la notación
\[
(a,b)\equiv a+bi.
\]

\textbf{Observación docente.}
Esta construcción muestra que los números complejos no son objetos
``misteriosos'': son pares ordenados de números reales equipados con una
multiplicación específicamente diseñada para que $(0,1)^2=(-1,0)$.

\paragraph{Ejemplo 1.}
En la construcción por pares,
\[
(2,3)+(4,-1)
=
(6,2).
\]
Esto corresponde a
\[
(2+3i)+(4-i)=6+2i.
\]

\paragraph{Ejemplo 2.}
\[
(1,2)(3,4)
=
(3-8,4+6)
=
(-5,10).
\]
En forma binómica:
\[
(1+2i)(3+4i)=-5+10i.
\]

\paragraph{Ejemplo 3.}
El real $7$ se identifica con
\[
(7,0).
\]
Entonces
\[
(7,0)+(0,2)=(7,2),
\]
es decir,
\[
7+2i.
\]

\paragraph{Ejemplo 4.}
La unidad imaginaria es
\[
i=(0,1).
\]
Por tanto,
\[
3i=(0,3).
\]

\paragraph{Ejemplo 5.}
El producto
\[
(0,2)(0,5)
\]
es
\[
(0\cdot0-2\cdot5,\;0\cdot5+2\cdot0)
=
(-10,0),
\]
por lo que
\[
(2i)(5i)=-10.
\]

% ============================================================
\subsection{Suma, resta y estructura aditiva}
% ============================================================

Sean
\[
z_1=a+bi,
\qquad
z_2=c+di.
\]

La suma se define componente a componente:
\[
\boxed{
z_1+z_2
=
(a+c)+(b+d)i.
}
\]

La resta se define mediante la suma con el inverso aditivo:
\[
z_1-z_2
=
z_1+(-z_2),
\]
de modo que
\[
\boxed{
z_1-z_2
=
(a-c)+(b-d)i.
}
\]

El inverso aditivo de
\[
z=a+bi
\]
es
\[
\boxed{-z=-a-bi.}
\]

El elemento neutro aditivo es
\[
0=0+0i.
\]

La suma de complejos satisface:
\[
z+w=w+z,
\]
\[
(z+w)+u=z+(w+u),
\]
\[
z+0=z,
\]
\[
z+(-z)=0.
\]

\paragraph{Ejemplo 1.}
\[
(3+2i)+(5-7i)=8-5i.
\]

\paragraph{Ejemplo 2.}
\[
(-4+6i)+(2+9i)=-2+15i.
\]

\paragraph{Ejemplo 3.}
\[
(7-3i)-(2+5i)
=
5-8i.
\]

\paragraph{Ejemplo 4.}
El inverso aditivo de
\[
z=4-9i
\]
es
\[
-z=-4+9i.
\]

\paragraph{Ejemplo 5.}
Simplifiquemos
\[
(2+3i)-(5-i)+(4+2i).
\]
Agrupando partes reales e imaginarias:
\[
(2-5+4)+(3+1+2)i
=
\boxed{1+6i}.
\]

% ============================================================
\subsection{Producto de números complejos}
% ============================================================

Sean
\[
z_1=a+bi,
\qquad
z_2=c+di.
\]

El producto se obtiene utilizando distributividad y la relación $i^2=-1$:
\begin{align*}
(a+bi)(c+di)
&=
ac+adi+bci+bdi^2\\
&=
ac+adi+bci-bd.
\end{align*}

Por tanto,
\[
\boxed{
(a+bi)(c+di)
=
(ac-bd)+(ad+bc)i.
}
\]

El producto es cerrado en $\mathbb{C}$, es decir,
\[
z,w\in\mathbb{C}
\quad\Longrightarrow\quad
zw\in\mathbb{C}.
\]

Además, satisface
\[
zw=wz,
\]
\[
(zw)u=z(wu),
\]
\[
z(w+u)=zw+zu.
\]

\paragraph{Ejemplo 1.}
\[
(2+3i)(4-i)
=
8-2i+12i-3i^2
=
\boxed{11+10i}.
\]

\paragraph{Ejemplo 2.}
\[
(1+i)^2
=
1+2i+i^2
=
\boxed{2i}.
\]

\paragraph{Ejemplo 3.}
\[
(3-2i)(3+2i)
=
9-4i^2
=
\boxed{13}.
\]

\paragraph{Ejemplo 4.}
\[
(-2+5i)(1-3i)
=
-2+6i+5i-15i^2
=
\boxed{13+11i}.
\]

\paragraph{Ejemplo 5.}
Simplifiquemos
\[
(1+i)(1-i)(2+i).
\]
Primero,
\[
(1+i)(1-i)=1-i^2=2.
\]
Por tanto,
\[
2(2+i)
=
\boxed{4+2i}.
\]

% ============================================================
\subsection{Conjugado complejo}
% ============================================================

\textbf{Definición.}
Sea
\[
z=a+bi.
\]
El \emph{conjugado complejo} de $z$ es
\[
\boxed{\overline z=a-bi.}
\]

La conjugación cambia el signo de la parte imaginaria y conserva la parte real.

\textbf{Proposición.}
Para $z,w\in\mathbb{C}$:
\[
\boxed{\overline{\overline z}=z,}
\]
\[
\boxed{\overline{z+w}=\overline z+\overline w,}
\]
\[
\boxed{\overline{zw}=\overline z\,\overline w.}
\]

Además,
\[
z+\overline z=2\operatorname{Re}(z),
\]
y
\[
z-\overline z=2i\operatorname{Im}(z).
\]

Por tanto,
\[
\boxed{
\operatorname{Re}(z)
=
\frac{z+\overline z}{2}
}
\]
y
\[
\boxed{
\operatorname{Im}(z)
=
\frac{z-\overline z}{2i}.
}
\]

\textbf{Caracterización.}
\[
z\in\mathbb{R}
\iff
z=\overline z.
\]

\paragraph{Ejemplo 1.}
Si
\[
z=3+5i,
\]
entonces
\[
\overline z=3-5i.
\]

\paragraph{Ejemplo 2.}
Si
\[
z=-4-7i,
\]
entonces
\[
\overline z=-4+7i.
\]

\paragraph{Ejemplo 3.}
Para
\[
z=2+3i,
\]
\[
z+\overline z
=
(2+3i)+(2-3i)
=
4,
\]
y por tanto
\[
\operatorname{Re}(z)=2.
\]

\paragraph{Ejemplo 4.}
Para
\[
z=1-4i,
\]
\[
z-\overline z
=
(1-4i)-(1+4i)
=
-8i.
\]
Entonces
\[
\operatorname{Im}(z)
=
\frac{-8i}{2i}
=
-4.
\]

\paragraph{Ejemplo 5.}
Verifiquemos
\[
\overline{zw}=\overline z\,\overline w
\]
para
\[
z=1+2i,
\qquad
w=3-i.
\]
Primero,
\[
zw=5+5i,
\]
por lo que
\[
\overline{zw}=5-5i.
\]
Además,
\[
\overline z=1-2i,
\qquad
\overline w=3+i,
\]
y
\[
(1-2i)(3+i)=5-5i.
\]
Por tanto, la identidad se verifica.

% ============================================================
\subsection{Módulo de un número complejo}
% ============================================================

\textbf{Definición.}
Para
\[
z=a+bi,
\]
se define el módulo de $z$ mediante
\[
\boxed{
|z|=\sqrt{a^2+b^2}.
}
\]

Como
\[
z\overline z
=
(a+bi)(a-bi)
=
a^2+b^2,
\]
se obtiene
\[
\boxed{
z\overline z=|z|^2.
}
\]

De esta identidad se deduce que
\[
|z|\geq 0
\]
y
\[
\boxed{|z|=0\iff z=0.}
\]

También
\[
|\overline z|=|z|.
\]

\textbf{Observación.}
El módulo será interpretado geométricamente como la distancia del punto
$z$ al origen en el plano complejo.

\paragraph{Ejemplo 1.}
Si
\[
z=3+4i,
\]
entonces
\[
|z|=\sqrt{3^2+4^2}=5.
\]

\paragraph{Ejemplo 2.}
Si
\[
z=-5+12i,
\]
entonces
\[
|z|=\sqrt{25+144}=13.
\]

\paragraph{Ejemplo 3.}
Si
\[
z=7i,
\]
entonces
\[
|z|=7.
\]

\paragraph{Ejemplo 4.}
Si
\[
z=-6,
\]
entonces
\[
|z|=6.
\]

\paragraph{Ejemplo 5.}
Sea
\[
z=2-3i.
\]
Entonces
\[
z\overline z
=
(2-3i)(2+3i)
=
4+9
=
13.
\]
Por otra parte,
\[
|z|^2
=
(\sqrt{13})^2
=
13.
\]
Así,
\[
z\overline z=|z|^2.
\]

% ============================================================
\subsection{División e inverso multiplicativo}
% ============================================================

Sea
\[
z=a+bi\neq0.
\]

Buscamos un número
\[
z^{-1}=x+yi
\]
tal que
\[
zz^{-1}=1.
\]

Utilizando el conjugado:
\[
z\overline z=|z|^2.
\]
Por tanto,
\[
z\frac{\overline z}{|z|^2}=1,
\]
y obtenemos
\[
\boxed{
z^{-1}
=
\frac{\overline z}{|z|^2},
\qquad z\neq0.
}
\]

Si
\[
z_1=a+bi
\]
y
\[
z_2=c+di\neq0,
\]
entonces
\[
\frac{z_1}{z_2}
=
z_1z_2^{-1}
=
\frac{z_1\overline{z_2}}{|z_2|^2}.
\]

Expresado en coordenadas:
\[
\boxed{
\frac{a+bi}{c+di}
=
\frac{ac+bd}{c^2+d^2}
+
\frac{bc-ad}{c^2+d^2}i.
}
\]

\paragraph{Ejemplo 1.}
Calculemos
\[
\frac{1}{2+i}.
\]
Multiplicando por el conjugado:
\[
\frac{1}{2+i}
\frac{2-i}{2-i}
=
\frac{2-i}{5}.
\]
Por tanto,
\[
\boxed{
\frac{1}{2+i}
=
\frac25-\frac15 i.
}
\]

\paragraph{Ejemplo 2.}
\[
\frac{3+2i}{1-i}
=
\frac{(3+2i)(1+i)}{(1-i)(1+i)}
=
\frac{1+5i}{2}.
\]
Así,
\[
\boxed{
\frac{3+2i}{1-i}
=
\frac12+\frac52 i.
}
\]

\paragraph{Ejemplo 3.}
El inverso de
\[
z=3-4i
\]
es
\[
z^{-1}
=
\frac{3+4i}{25}
=
\boxed{
\frac3{25}+\frac4{25}i
}.
\]

\paragraph{Ejemplo 4.}
\[
\frac{2-i}{2+i}
=
\frac{(2-i)^2}{5}
=
\frac{3-4i}{5}.
\]
Por tanto,
\[
\boxed{
\frac{2-i}{2+i}
=
\frac35-\frac45 i.
}
\]

\paragraph{Ejemplo 5.}
Simplifiquemos
\[
\frac{i}{1+i}.
\]
Entonces
\[
\frac{i}{1+i}
\frac{1-i}{1-i}
=
\frac{i-i^2}{2}
=
\frac{1+i}{2}.
\]
Por tanto,
\[
\boxed{
\frac{i}{1+i}
=
\frac12+\frac12 i.
}
\]

% ============================================================
\subsection{$\mathbb{C}$ como campo}
% ============================================================

El conjunto $\mathbb{C}$, equipado con las operaciones usuales de suma y
producto, constituye un campo.

Esto significa que para cualesquiera $z,w,u\in\mathbb{C}$ se satisfacen
los axiomas:

\begin{enumerate}
    \item cerradura de la suma;
    \item asociatividad de la suma;
    \item conmutatividad de la suma;
    \item existencia del neutro aditivo $0$;
    \item existencia del inverso aditivo $-z$;
    \item cerradura del producto;
    \item asociatividad del producto;
    \item conmutatividad del producto;
    \item existencia del neutro multiplicativo $1$;
    \item existencia del inverso multiplicativo $z^{-1}$ para $z\neq0$;
    \item distributividad del producto respecto de la suma.
\end{enumerate}

Por tanto,
\[
\boxed{\mathbb{C}\text{ es un campo}.}
\]

Además,
\[
\boxed{\mathbb{R}\subset\mathbb{C}}
\]
y las operaciones de $\mathbb{C}$ restringidas a $\mathbb{R}$ coinciden con
las operaciones usuales de los números reales.

\textbf{Teorema.}
No existe un orden total sobre $\mathbb{C}$ compatible con las operaciones
de campo que extienda el orden usual de $\mathbb{R}$.

\textbf{Idea de demostración.}
En un campo ordenado, el cuadrado de todo elemento no nulo debe ser positivo.
Si $\mathbb{C}$ fuera ordenado, entonces
\[
i^2>0.
\]
Pero
\[
i^2=-1,
\]
lo cual implicaría
\[
-1>0,
\]
contradicción con el orden usual de $\mathbb{R}$.
Por tanto, $\mathbb{C}$ no puede ser un campo ordenado.
\hfill$\square$

\paragraph{Ejemplo 1.}
La suma
\[
(2+i)+(3-5i)=5-4i
\]
permanece en $\mathbb{C}$, mostrando cerradura aditiva.

\paragraph{Ejemplo 2.}
El producto
\[
(1+i)(2-3i)=5-i
\]
permanece en $\mathbb{C}$, mostrando cerradura multiplicativa.

\paragraph{Ejemplo 3.}
El inverso aditivo de
\[
z=-2+7i
\]
es
\[
2-7i.
\]

\paragraph{Ejemplo 4.}
Para
\[
z=1+i\neq0,
\]
el inverso multiplicativo es
\[
z^{-1}
=
\frac{1-i}{2}.
\]
En efecto,
\[
(1+i)\frac{1-i}{2}
=
1.
\]

\paragraph{Ejemplo 5.}
La expresión
\[
i>0
\]
no tiene un significado compatible con la estructura de campo de
$\mathbb{C}$ como sí ocurre en $\mathbb{R}$. Esto explica por qué las
desigualdades complejas requieren conceptos diferentes, típicamente
basados en el módulo.

% ============================================================
\subsection{Resumen algebraico de la primera parte}
% ============================================================

La construcción algebraica de los números complejos puede resumirse mediante

\[
\boxed{
\begin{array}{c}
\mathbb{R}
\\[0.2cm]
\downarrow
\\[0.2cm]
i^2=-1
\\[0.2cm]
\downarrow
\\[0.2cm]
z=a+bi
\\[0.2cm]
\downarrow
\\[0.2cm]
\operatorname{Re}(z),\operatorname{Im}(z)
\\[0.2cm]
\downarrow
\\[0.2cm]
\text{suma y producto}
\\[0.2cm]
\downarrow
\\[0.2cm]
\overline z,\ |z|
\\[0.2cm]
\downarrow
\\[0.2cm]
z^{-1}=\dfrac{\overline z}{|z|^2}
\\[0.2cm]
\downarrow
\\[0.2cm]
\mathbb{C}\text{ es un campo}
\end{array}
}
\]

Esta primera parte proporciona la estructura algebraica necesaria para
desarrollar posteriormente la geometría del plano complejo, la forma polar,
la fórmula de Euler, el teorema de De Moivre, las raíces complejas y el
logaritmo complejo.
